% Part: incompleteness
% Chapter: theories-computability
% Section: extensions-of-q-not-decidable

\documentclass[../../../include/open-logic-section]{subfiles}

\begin{document}

\olfileid{inc}{tcp}{cqn}

\olsection{گسترش‌های سازگارِ $\Th{Q}$ تصمیم‌ناپذیرند}

\begin{explain}
به یاد آورید نظریه‌ای را \emph{سازگار} می‌نامیم اگر هم~$!A$ و هم
$\lnot !A$ را برای هیچ فرمول~$!A$ای ثابت نکند. چون هر چیزی از
یک تناقض نتیجه می‌شود، نظریهٔ ناسازگار بدیهی است: هر جمله‌ای در آن
اثبات‌پذیر است. روشن است که اگر نظریه‌ای $\omega$-سازگار باشد، سازگار
نیز هست. اما سازگاری شرط ضعیف‌تری است (یعنی نظریه‌هایی وجود دارند که
سازگارند اما $\omega$-سازگار نیستند). می‌توانیم فرض
\olref[oqn]{thm:oconsis-q} را به سازگاری ساده تضعیف کنیم و قضیه‌ای
قوی‌تر به دست آوریم.
\end{explain}

\begin{lem}
هیچ «رابطهٔ محاسبه‌پذیر جهان‌شمولی» وجود ندارد. یعنی هیچ رابطهٔ
محاسبه‌پذیر دوتایی~$R(x,y)$ای با ویژگی زیر وجود ندارد: هرگاه~$S(y)$
رابطه‌ای محاسبه‌پذیر و یک‌موضعی باشد، $k$ای وجود داشته باشد به‌گونه‌ای
که برای هر~$y$، $S(y)$ صادق باشد اگر و تنها اگر~$R(k,y)$ صادق باشد.
\end{lem}

\begin{proof}
فرض کنید~$R(x,y)$ رابطه‌ای محاسبه‌پذیر و جهان‌شمول باشد. بگذارید
$S(y)$ رابطهٔ~$\lnot R(y,y)$ باشد. چون~$S(y)$ محاسبه‌پذیر است، برای
بعضی~$k$، $S(y)$ با~$R(k,y)$ هم‌ارز است. اما در آن صورت~$S(k)$ هم
با~$R(k,k)$ و هم با~$\lnot R(k,k)$ هم‌ارز است، که تناقض است.
\end{proof}


\begin{thm}
بگذارید~$\Th{T}$ هر نظریهٔ سازگاری باشد که~$\Th{Q}$ را در بر
می‌گیرد. آنگاه~$\Th{T}$ تصمیم‌پذیر نیست.
\end{thm}


\begin{proof}
فرض کنید~$\Th{T}$ گسترشی سازگار و تصمیم‌پذیر از~$\Th{Q}$ باشد. با
استفاده از~$\Th{T}$ برای تعریف یک رابطهٔ محاسبه‌پذیر جهان‌شمول، به
تناقض خواهیم رسید.

بگذارید~$R(x,y)$ برقرار باشد اگر و تنها اگر
\begin{quote}
$x$ فرمولی چون~$!D(u)$ را رمزگذاری کند و~$\Th{T}$،
$!D(\num y)$ را ثابت کند.
\end{quote}
چون فرض کرده‌ایم~$\Th{T}$ تصمیم‌پذیر است، $R$ محاسبه‌پذیر است. نشان
می‌دهیم~$R$ جهان‌شمول است. اگر~$S(y)$ هر رابطهٔ محاسبه‌پذیری باشد،
آنگاه در~$\Th{Q}$ (و بنابراین در~$\Th{T}$) با فرمول~$!D_S(u)$
بازنمایی می‌شود. پس برای هر~$n$ داریم
\begin{eqnarray*}
S(n) & \lif & T \vdash !D_S(\num n) \\
& \lif & R(\Gn{!D_S(u)}, n)
\end{eqnarray*}
و
\begin{eqnarray*}
\lnot S(n) & \lif & T \vdash \lnot !D_S(\num n) \\
& \lif & T \not\vdash !D_S(\num n) \quad \text{(زیرا $\Th{T}$
  سازگار است)} \\
& \lif & \lnot R(\Gn{!D_S(u)}, n).
\end{eqnarray*}
یعنی برای هر~$y$، $S(y)$ صادق است اگر و تنها اگر
$R(\Gn{!D_S(u)}, y)$ صادق باشد. پس~$R$ جهان‌شمول است و به تناقضی که
در پی آن بودیم رسیده‌ایم.
\end{proof}

«حساب صادق» را نظریهٔ~$\Setabs{!A}{ \Sat{\Nat}{!A}}$ می‌نامیم؛ یعنی
مجموعهٔ جمله‌های زبان حساب که در تعبیر استاندارد صادق‌اند.

\begin{cor}
حساب صادق تصمیم‌پذیر نیست.
\end{cor}

%In Section~\olref{undefinability:truth:section} we will state a stronger
%result, due to Tarski.

\end{document}
